\documentclass[11pt]{article}

\usepackage[margin=2.3cm]{geometry}
\usepackage[T1]{fontenc}
\usepackage{lmodern}
\usepackage{amsmath,amssymb}
\usepackage{booktabs}
\usepackage{multirow}
\usepackage{caption}
\usepackage[colorlinks=true,linkcolor=blue!50!black]{hyperref}

\captionsetup{font=small,labelfont=bf,justification=justified,singlelinecheck=false}
\newcommand{\code}[1]{\texttt{#1}}
\newcommand{\best}[1]{\textbf{#1}}
\newcommand{\ci}[2]{\,[#1,\,#2]}

\title{ASD-ContextPose: Results Tables}
\author{ASD-ContextPose project}
\date{8 November 2026}

\begin{document}
\maketitle

%=====================================================================
% Binary calls (phase 2): main results
%=====================================================================

\begin{table}[!htbp]
\centering
\caption{\textbf{Stream ablation of ContextPose on the binary valence and arousal calls.}
\emph{Model:} ContextPose (context-queried temporal attention over the pre-event, event and post-event windows, with an optional gated face stream), with each input stream used alone, in pairs, in threes and all together.
\emph{Streams:} P = upper-body pose (11 joints, $x$, $y$, confidence and real-time speed); A = event annotation (four binary context indicators); V = VLM scene description (Qwen3-Omni window descriptions embedded with MiniLM, 384-d), used as the context vector in place of or alongside A; F = facial expression (frozen FER ViT embedding per step, kept only where the face passes a frontal/YuNet/brightness quality gate).
\emph{Offsets:} rows with A start from the event-annotation lookup's log-odds and learn a correction to it; rows without A start from a constant offset fitted on the training folds, so no annotation information reaches them.
\emph{Data and protocol:} 153 development clips from 68 videos; video-grouped 5-fold cross-validation (duplicate videos merged); each clip predicted once by a model trained on three folds and early-stopped on the fourth. All rows share the folds, the votes objective (binary cross-entropy on annotator vote shares) and the same training settings (width 64, dropout 0.3, weight decay 0.05, learning rate $3\times10^{-4}$, pose augmentation, up to 60 epochs with patience 10).
\emph{Targets:} annotator-majority calls, ties excluded: valence positive vs not positive (150 clips, 100 positive) and arousal high vs low (148 clips, 33 high).
\emph{Metrics:} 5-seed ensembles (majority of seed calls; mean seed score for AUC). Acc.\ = accuracy; Bal.\ acc.\ = balanced accuracy (mean per-class recall, chance 50\%); AUC = area under the ROC curve. Decision thresholds maximise balanced accuracy on each model's own training and early-stopping clips, never on the scored fold. Brackets are 95\% participant-bootstrap intervals (2000 draws).
\emph{References:} annotators are each annotator scored against the other two where those two agree; the lookup predicts the training folds' mean vote share for each indicator combination. Bold: best model in each column.}
\label{tab:streams}
\small
\setlength{\tabcolsep}{3.5pt}
\begin{tabular}{cccclcccccc}
\toprule
 & & & & & \multicolumn{3}{c}{Valence: positive vs not ($n{=}150$)} & \multicolumn{3}{c}{Arousal: high vs low ($n{=}148$)} \\
\cmidrule(lr){6-8}\cmidrule(lr){9-11}
P & A & V & F & Model & Acc. & Bal. acc. & AUC & Acc. & Bal. acc. & AUC \\
\midrule
\multicolumn{5}{l}{Annotators (each vs the other two)} & 79\% & 79\% & -- & 81\% & 77\% & -- \\
\multicolumn{5}{l}{Always majority} & 67\% & 50\% & 0.50 & 78\% & 50\% & 0.50 \\
 & $\checkmark$ & & & Event-annotation lookup (no network) & 55\% & 53\% & 0.55 & 60\% & 68\% & 0.63 \\
\midrule
$\checkmark$ & & & & Pose only & 48\% & 46\%\ci{34}{59} & 0.52 & 55\% & 53\%\ci{40}{68} & 0.57 \\
 & $\checkmark$ & & & Event annotation only & 58\% & 55\%\ci{42}{68} & 0.58 & 60\% & 68\%\ci{55}{79} & 0.61 \\
 & & $\checkmark$ & & VLM only & 49\% & 48\%\ci{36}{61} & 0.53 & 76\% & \best{76\%}\ci{66}{86} & 0.69 \\
 & & & $\checkmark$ & Face only & 60\% & 60\%\ci{46}{72} & 0.73 & 57\% & 65\%\ci{54}{74} & 0.69 \\
\midrule
$\checkmark$ & $\checkmark$ & & & Pose + annotation & 53\% & 54\%\ci{42}{63} & 0.60 & 61\% & 62\%\ci{50}{74} & 0.68 \\
$\checkmark$ & & $\checkmark$ & & Pose + VLM & 49\% & 48\%\ci{35}{61} & 0.52 & 60\% & 60\%\ci{46}{71} & 0.55 \\
$\checkmark$ & & & $\checkmark$ & Pose + face & 49\% & 50\%\ci{36}{61} & 0.58 & 59\% & 55\%\ci{41}{68} & 0.61 \\
 & $\checkmark$ & $\checkmark$ & & Annotation + VLM & 63\% & 61\%\ci{50}{70} & 0.60 & 75\% & 69\%\ci{60}{81} & \best{0.75} \\
 & $\checkmark$ & & $\checkmark$ & Annotation + face (V3) & 67\% & \best{65\%}\ci{53}{75} & 0.71 & 68\% & 69\%\ci{57}{79} & 0.73 \\
 & & $\checkmark$ & $\checkmark$ & VLM + face & 62\% & 62\%\ci{49}{74} & \best{0.74} & 70\% & 69\%\ci{56}{79} & 0.72 \\
\midrule
$\checkmark$ & $\checkmark$ & & $\checkmark$ & \textbf{Pose + annotation + face (full model)} & 59\% & 58\%\ci{46}{69} & 0.63 & 61\% & 60\%\ci{48}{74} & 0.69 \\
$\checkmark$ & & $\checkmark$ & $\checkmark$ & Pose + VLM + face (VLM in place of annotation) & 51\% & 49\%\ci{35}{63} & 0.58 & 64\% & 64\%\ci{52}{75} & 0.64 \\
$\checkmark$ & $\checkmark$ & $\checkmark$ & & Pose + annotation + VLM & 57\% & 58\%\ci{46}{67} & 0.59 & 61\% & 59\%\ci{47}{71} & 0.67 \\
 & $\checkmark$ & $\checkmark$ & $\checkmark$ & Annotation + VLM + face & 65\% & 64\%\ci{52}{75} & 0.70 & 73\% & 69\%\ci{56}{79} & \best{0.75} \\
\midrule
$\checkmark$ & $\checkmark$ & $\checkmark$ & $\checkmark$ & All four streams & 60\% & 58\%\ci{44}{70} & 0.66 & 58\% & 56\%\ci{43}{70} & 0.70 \\
\bottomrule
\end{tabular}
\end{table}

\begin{table}[!htbp]
\centering
\caption{\textbf{Best cross-validated results per target, across every feature set tried.}
Binary calls from the annotators' majority (ties excluded): valence positive vs not positive (150 clips, 100 positive) and arousal high vs low (148 clips, 33 high), on 153 development clips with video-grouped 5-fold cross-validation.
\emph{Predictors:} V3 = ContextPose starting from the event-annotation lookup's log-odds plus a learned facial-expression correction (no pose). Lookup + VLM = ContextPose on the VLM scene-description embedding, starting from the lookup, trained for up to 150 epochs. VLM embedding only = the same network with a constant offset instead of the lookup, so it uses no annotation (up to 60 epochs).
\emph{Metrics:} 5-seed ensembles trained on annotator vote shares; balanced accuracy is the mean per-class recall (chance 50\%); thresholds are set on each model's own training clips. Brackets are 95\% participant-bootstrap intervals.
\emph{References:} annotators are each annotator scored against the other two where those two agree; always majority predicts the commoner class (positive valence, low arousal). Bold: best balanced accuracy and AUC per target. No gain over the lookup has an interval that excludes zero.}
\label{tab:headline}
\begin{tabular}{llccc}
\toprule
Call & Predictor & Accuracy & Balanced accuracy & AUC \\
\midrule
\multirow{4}{*}{\shortstack[l]{Valence:\\positive vs not}}
 & Annotators (each vs the other two) & 79\% & 79\% & -- \\
 & Always majority & 67\% & 50\% & 0.50 \\
 & Event-context lookup & 55\% & 53\% & 0.55 \\
 & \best{V3} (context lookup + face) & 67\% & \best{65\%}\ci{53}{75} & \best{0.71}\ci{0.56}{0.83} \\
\midrule
\multirow{6}{*}{\shortstack[l]{Arousal:\\high vs low}}
 & Annotators (each vs the other two) & 81\% & 77\% & -- \\
 & Always majority & 78\% & 50\% & 0.50 \\
 & Event-context lookup & 60\% & 68\% & 0.63 \\
 & V3 & 68\% & 69\%\ci{57}{79} & 0.73\ci{0.62}{0.85} \\
 & \best{Lookup + VLM embedding} (150-epoch cap) & 79\% & 72\%\ci{63}{85} & \best{0.77}\ci{0.66}{0.87} \\
 & \best{VLM embedding only}, no annotation & 76\% & \best{76\%}\ci{66}{86} & 0.69\ci{0.57}{0.83} \\
\bottomrule
\end{tabular}
\end{table}

\clearpage
%=====================================================================
% Data
%=====================================================================

\begin{table}[!htbp]
\centering
\caption{\textbf{Missing child pose before and after gap filling} (all 192 context clips, 159 train and 33 test).
\code{ViTPoseClean} is the upstream ViTPose pipeline (11 upper-body COCO joints at 25\,fps). Gap fill pass 1 recovers untracked frames from the source videos with RT-DETR person detection and ViTPose-base, accepting the person closest to the previous pose and stopping after 0.5\,s without a match. Pass 2 re-acquires the child after long gaps by the colour of their clothing.
\emph{Columns:} frames without pose are counted over all clips; an empty window is a pre-event, event or post-event window with no tracked pose at all; usable clips have no empty window and a non-empty event (one 10\,s clip has no event frames). Bold: final number of usable clips.}
\label{tab:gapfill}
\begin{tabular}{lccc}
\toprule
Stage & Frames without pose & Clips with an empty window & Usable clips \\
\midrule
\code{ViTPoseClean} & 21{,}026 (19\%) & 43 (22.4\%) & 148 \\
Gap fill, pass 1 & 10{,}815 (10\%) & 9 (4.7\%) & 182 \\
Pass 2 (re-acquisition) & -- & 6 (3.1\%) & \best{185} \\
\bottomrule
\end{tabular}
\end{table}

\begin{table}[!htbp]
\centering
\caption{\textbf{The event-annotation lookup}, fitted on all 153 development clips (in cross-validation it is refitted on each fold's training clips).
Each row is one combination of the four binary context indicators that annotators coded for the event (social interaction, auditory event, activity or object change, visual or media event); a blank row of indicators means no event. Valence and arousal are the median consensus labels of the clips with that combination, on a $[-1, 1]$ scale (positive = positive valence or high arousal). Combinations with fewer than 3 clips use the overall medians. Only 9 of the 16 possible combinations occur.}
\label{tab:lookup}
\begin{tabular}{ccccrrr}
\toprule
Social & Auditory & Object change & Visual/media & Clips & Valence & Arousal \\
\midrule
 & $\checkmark$ & & & 57 & $+0.69$ & $-0.38$ \\
 & $\checkmark$ & & $\checkmark$ & 25 & $0.00$ & $-1.00$ \\
$\checkmark$ & $\checkmark$ & & & 18 & $+0.70$ & $-0.41$ \\
$\checkmark$ & $\checkmark$ & $\checkmark$ & & 15 & $+1.00$ & $-0.43$ \\
$\checkmark$ & $\checkmark$ & & $\checkmark$ & 11 & $+1.00$ & $-0.45$ \\
$\checkmark$ & & & & 11 & $+0.67$ & $-1.00$ \\
 & & & & 10 & $0.00$ & $-1.00$ \\
 & $\checkmark$ & $\checkmark$ & & 4 & $+0.82$ & $-0.15$ \\
 & $\checkmark$ & $\checkmark$ & $\checkmark$ & 2 & \multicolumn{2}{c}{overall: $+0.69$ / $-0.64$} \\
\bottomrule
\end{tabular}
\end{table}

\clearpage
%=====================================================================
% Phase 1: regression
%=====================================================================

\begin{table}[!htbp]
\centering
\caption{\textbf{Phase 1 architecture comparison: regression of the consensus valence and arousal labels} ($[-1, 1]$ scale) on 153 development clips, video-grouped 5-fold cross-validation (original folds, before the duplicate-video merge), predictions averaged over 5 seeds.
\emph{Models:} \code{text\_attention} = ContextPose conditioned on the CLIP ViT-L/14 embedding of the event sentence; \code{annotation\_attention} = the same conditioned on the four indicators; \code{annotation\_film} = indicators modulating the pose features by FiLM; pose only = context zeroed.
\emph{Training:} default = width 128, learning rate $10^{-3}$, 30 epochs, patience 5; regularised = width 64, dropout 0.3, weight decay 0.05, learning rate $3\times10^{-4}$, pose augmentation.
\emph{Metrics:} MAE = mean absolute error (lower is better); CCC = concordance correlation coefficient (higher is better; 0 for a constant prediction). Mean MAE averages valence (V) and arousal (A). Bold: best, the non-learned lookup, which no network beats.}
\label{tab:spec}
\begin{tabular}{llcccc}
\toprule
Model & Training & Mean MAE & V MAE / CCC & A MAE / CCC \\
\midrule
\code{text\_attention} (CLIP) & default & 0.483 & 0.400 / $-0.06$ & 0.567 / $-0.04$ \\
\code{annotation\_attention} & default & 0.482 & 0.409 / $-0.04$ & 0.554 / $0.02$ \\
\code{annotation\_film} & default & 0.496 & 0.421 / $-0.08$ & 0.572 / $0.03$ \\
Pose only & default & 0.496 & 0.428 / $-0.12$ & 0.563 / $0.01$ \\
\code{text\_attention} (CLIP) & regularised & 0.467 & 0.383 / $-0.03$ & 0.551 / $-0.01$ \\
\code{annotation\_attention} & regularised & 0.468 & 0.383 / $-0.01$ & 0.553 / $-0.03$ \\
\code{annotation\_film} & regularised & 0.468 & 0.389 / $0.03$ & 0.547 / $0.07$ \\
Pose only & regularised & 0.468 & 0.384 / $-0.01$ & 0.553 / $-0.03$ \\
\midrule
Overall median & -- & 0.464 & 0.373 / $0.00$ & 0.555 / $-0.08$ \\
\best{Context lookup (median per combination)} & -- & \best{0.425} & \best{0.351} / \best{0.34} & \best{0.498} / \best{0.15} \\
\bottomrule
\end{tabular}
\end{table}

\begin{table}[!htbp]
\centering
\caption{\textbf{Phase 1: residual learning, L1 loss and a real-time speed channel} (regression of the consensus labels, 153 development clips, video-grouped 5-fold cross-validation on the original folds, predictions averaged over 5 seeds).
\emph{Configurations:} residual = the network predicts a correction to the training folds' lookup median; speed = a fourth pose feature with joint speed in shoulder widths per second at 25\,fps, which keeps real-time units after each window is resampled to 32 steps; L1 / MSE = training loss; regularisation as in Table~\ref{tab:spec}; S5 = S4 with the context removed from the network (the residual still supplies the lookup).
\emph{Metrics:} MAE (lower is better) and CCC (higher is better) for valence (V) and arousal (A). Gain = lookup mean MAE minus model mean MAE (positive = better than the lookup), with a 95\% participant-bootstrap interval where computed. Bold: best, the lookup.}
\label{tab:speed}
\begin{tabular}{lcccc}
\toprule
Configuration & Mean MAE & V MAE / CCC & A MAE / CCC & Gain vs lookup \\
\midrule
\best{Context lookup} & \best{0.425} & 0.351 / 0.34 & 0.498 / 0.15 & -- \\
S6: residual + speed, L1 & 0.434 & 0.359 / 0.36 & 0.510 / 0.19 & $-0.009$\ci{-0.049}{+0.028} \\
S4: residual + speed, MSE & 0.436 & 0.362 / 0.38 & 0.511 / 0.25 & $-0.012$\ci{-0.054}{+0.033} \\
S5: as S4, pose only & 0.436 & 0.362 / 0.38 & 0.509 / 0.25 & $-0.011$\ci{-0.054}{+0.033} \\
Residual, no speed (MSE) & 0.451 & -- & -- & -- \\
Residual, no speed (L1) & 0.445 & 0.373 / 0.27 & 0.517 / 0.15 & -- \\
S2: speed + regularisation & 0.445 & 0.380 / 0.09 & 0.509 / 0.10 & $-0.020$ \\
S1: speed, no regularisation & 0.459 & 0.411 / 0.01 & 0.507 / 0.17 & $-0.034$ \\
\bottomrule
\end{tabular}
\end{table}

\begin{table}[!htbp]
\centering
\caption{\textbf{Held-out test set, scored once} (29 usable test clips from 20 videos not seen in development; regression of the consensus labels). All predictors were fitted on the 153 development clips, and the candidate was fixed in advance from cross-validation: the lookup for valence and a linear model on context plus nine hand-crafted movement features per window (head, elbow, wrist and body speed, wrist height and spread, head bob) for arousal. The movement signal did not replicate (arousal CCC 0.00); the lookup's cross-validated performance carried over, making it the only predictor confirmed on new videos. MAE lower is better; CCC higher is better (0.00 = no agreement). Bold: best.}
\label{tab:test}
\begin{tabular}{lccc}
\toprule
Predictor & Mean MAE & V MAE / CCC & A MAE / CCC \\
\midrule
Candidate (lookup valence + linear movement arousal) & 0.428 & 0.361 / 0.31 & 0.495 / 0.00 \\
\best{Context lookup} & \best{0.389} & 0.361 / 0.31 & \best{0.417} / \best{0.17} \\
Overall median & 0.456 & 0.426 / 0.00 & 0.486 / 0.00 \\
Linear context + movement (both targets) & 0.450 & 0.406 / 0.12 & 0.495 / 0.00 \\
\bottomrule
\end{tabular}
\end{table}

\clearpage
%=====================================================================
% Face, audio and video features
%=====================================================================

\begin{table}[!htbp]
\centering
\caption{\textbf{Frozen face, video and audio features as a linear correction on the lookup} (regression of the consensus labels, 153 development clips, video-grouped 5-fold cross-validation on the original folds). Each embedding feeds an L1-penalised linear model that corrects the training folds' lookup.
\emph{Features:} face = crops at the tracked face keypoints embedded with the FER ViT \code{trpakov/vit-face-expression} (768-d [CLS]) or its 7 expression probabilities; video = V-JEPA 2 on a child crop; audio = AST embedding, or AudioSet event probabilities plus loudness.
\emph{Metrics:} MAE (lower is better) and CCC (higher is better) for valence (V) and arousal (A). Gain = lookup mean MAE minus model mean MAE (positive = better), with a 95\% participant-bootstrap interval for the face rows. Bold: best.}
\label{tab:media}
\begin{tabular}{lcccc}
\toprule
Input & Mean MAE & V MAE / CCC & A MAE / CCC & Gain vs lookup \\
\midrule
Context lookup & 0.425 & 0.351 / 0.34 & 0.498 / 0.15 & -- \\
\best{Face, FER embedding} & \best{0.420} & \best{0.328} / \best{0.47} & 0.511 / 0.21 & $+0.005$\ci{-0.014}{+0.026} \\
Face, 7 expression probabilities & 0.421 & 0.330 / 0.41 & 0.512 / 0.16 & $+0.004$\ci{-0.024}{+0.026} \\
Video, V-JEPA 2 (child crop) & 0.426 & 0.351 / 0.34 & 0.501 / 0.15 & $-0.001$ \\
Audio, AudioSet probabilities + loudness & 0.429 & 0.357 / 0.32 & 0.502 / 0.14 & $-0.005$ \\
Audio, AST embedding & 0.437 & 0.353 / 0.34 & 0.520 / 0.15 & $-0.012$ \\
\bottomrule
\end{tabular}
\end{table}

\begin{table}[!htbp]
\centering
\caption{\textbf{Phase 1: the face stream inside ContextPose} (regression of the consensus labels, 153 development clips, video-grouped 5-fold cross-validation on the original folds, predictions averaged over 5 seeds). The face stream encodes the FER embedding at each of the 32 steps per window, averages the steps that pass the quality gate (frontal face, YuNet detector score $\geq 0.8$, brightness $\geq 40$, sharpness $\geq 10$; 39\% of steps kept, 27 clips with no usable face), and adds a zero-initialised correction to the prediction.
\emph{Configurations:} F0 = residual on the lookup + pose with speed, no face; F1 = F0 + face correcting valence only; F2 = F0 + face correcting both targets; F3 = lookup + face correction, no pose.
\emph{Metrics:} MAE (lower is better) and valence CCC (higher is better). Gain = mean MAE of the training-fold lookup these models start from (0.422 mean; 0.368 valence, 0.476 arousal) minus the model's, with a 95\% participant-bootstrap interval for F3. Bold: best.}
\label{tab:face}
\begin{tabular}{lccccc}
\toprule
Configuration & Mean MAE & V MAE & V CCC & A MAE & Mean gain vs lookup \\
\midrule
F0: residual + speed, no face & 0.434 & 0.359 & 0.36 & 0.510 & $-0.012$ \\
F1: + face $\to$ valence & 0.430 & 0.346 & 0.45 & 0.513 & $-0.008$ \\
F2: + face $\to$ both & 0.425 & 0.347 & 0.43 & 0.503 & $-0.003$ \\
\best{F3: lookup + face, no pose} & \best{0.411} & 0.347 & 0.44 & 0.476 & $+0.011$\ci{-0.011}{+0.029} \\
\bottomrule
\end{tabular}
\end{table}

\clearpage
%=====================================================================
% Phase 2: objectives, EmotionCLIP, VLM, recurrent baselines
%=====================================================================

\begin{table}[!htbp]
\centering
\caption{\textbf{Regression vs votes objective on the binary calls} (153 development clips, video-grouped 5-fold cross-validation with duplicate videos merged, 5-seed ensembles). All models are \code{annotation\_attention} ContextPose; thresholds maximise balanced accuracy on each model's own training clips.
\emph{Objectives:} R = trained to regress the confidence-weighted consensus label; V = trained with binary cross-entropy on the share of annotators voting positive valence / high arousal.
\emph{Inputs:} pose + context = pose with speed and the four indicators, starting from the lookup; context + face = lookup plus a face correction, no pose (V3 is the votes version); face only = face correction on a constant offset.
\emph{Targets:} annotator-majority calls, ties excluded: valence positive vs not (150 clips, 100 positive), arousal high vs low (148 clips, 33 high). Balanced accuracy is the mean per-class recall (chance 50\%). Bold: best per column.}
\label{tab:v3}
\small
\begin{tabular}{lcccccc}
\toprule
 & \multicolumn{3}{c}{Valence: positive vs not} & \multicolumn{3}{c}{Arousal: high vs low} \\
\cmidrule(lr){2-4}\cmidrule(lr){5-7}
Predictor & Acc. & Bal. acc. & AUC & Acc. & Bal. acc. & AUC \\
\midrule
Annotators (ceiling) & 79\% & 79\% & -- & 81\% & 77\% & -- \\
Always majority & 67\% & 50\% & 0.50 & 78\% & 50\% & 0.50 \\
Context lookup & 55\% & 53\% & 0.55 & 60\% & 68\% & 0.63 \\
R: pose + context & 62\% & 60\% & 0.58 & 66\% & 63\% & 0.66 \\
R: context + face & 67\% & 62\% & 0.60 & 74\% & 69\% & 0.66 \\
R: face only & 59\% & 59\% & 0.70 & 64\% & 62\% & 0.67 \\
V: pose + context & 53\% & 54\% & 0.60 & 61\% & 62\% & 0.68 \\
\best{V3: context + face} & 67\% & \best{65\%} & 0.71 & 68\% & \best{69\%} & \best{0.73} \\
V: face only & 60\% & 60\% & \best{0.73} & 57\% & 65\% & 0.69 \\
\bottomrule
\end{tabular}
\end{table}

\begin{table}[!htbp]
\centering
\caption{\textbf{EmotionCLIP features on the binary calls} (153 development clips, video-grouped 5-fold cross-validation with duplicate videos merged). EmotionCLIP (ViT-B/32 with a subject mask and a temporal transformer) was used frozen, with the child's pose-derived box as the subject mask: a 512-d embedding per step and per window.
\emph{Approaches:} linear = class-balanced logistic regression on the one-hot indicator combination plus PCA-reduced embeddings (8 or 16 components, fitted in-fold), with the regularisation chosen on an inner fold; network = ContextPose with an EmotionCLIP stream in place of the face stream, starting from the lookup (V3 settings, 5-seed ensemble).
\emph{Ranges:} for the linear rows, the value range spans the PCA sizes, averaged over 5 fold arrangements; the parenthesised range is the spread over those arrangements. Brackets on network rows are 95\% participant-bootstrap intervals. Targets: arousal high vs low (148 clips, 33 high); valence positive vs not (150 clips, 100 positive). Bold: best.}
\label{tab:eclip}
\begin{tabular}{llcc}
\toprule
Call & Approach & Balanced accuracy & AUC \\
\midrule
\multirow{4}{*}{Arousal}
 & Context lookup & 68\% & 0.63 \\
 & EmotionCLIP alone (linear) & -- & 0.72--0.77 \\
 & \best{Linear: context + face + EmotionCLIP} & \best{74--75\%} (65--79\%) & \best{0.80--0.81} \\
 & Network: context + EmotionCLIP stream & 60\%\ci{47}{75} & 0.75 \\
\midrule
\multirow{2}{*}{Valence}
 & EmotionCLIP alone (linear) & 41--47\% & 0.35--0.48 \\
 & Network: context + EmotionCLIP stream & 52\%\ci{42}{61} & 0.59 \\
\bottomrule
\end{tabular}
\end{table}

\begin{table}[!htbp]
\centering
\caption{\textbf{VLM scene descriptions and a second (uploaded) face model on the binary calls} (153 development clips, video-grouped 5-fold cross-validation with duplicate videos merged, votes objective, 5-seed ensembles, V3 training settings). Cells are balanced accuracy / AUC.
\emph{Features:} VLM = Qwen3-Omni descriptions of 5--6\,s windows (scene, child actions, audible context) embedded with \code{all-MiniLM-L6-v2} and averaged into one 384-d clip vector, used as ContextPose's context vector; our FER face = the detector-gated FER ViT embedding; uploaded ViT face = \code{face\_features\_v3/vit} \code{backbone\_cls} at the frame nearest each step, masked by its own face detection (cosine similarity with our FER embedding 0.17).
\emph{Offsets:} every row above the rule starts from the event-annotation lookup, so ``lookup + VLM'' combines the annotation and the VLM; the last row uses a constant offset and so replaces the annotation with the VLM. Rows are trained for up to 60 epochs unless marked 150. Valence: positive vs not (150 clips, 100 positive); arousal: high vs low (148 clips, 33 high). Bold: best per column.}
\label{tab:vlm}
\begin{tabular}{lcc}
\toprule
Model & Valence & Arousal \\
\midrule
V3: lookup + our FER face & \best{65\% / 0.71} & 69\% / 0.73 \\
F3: lookup + uploaded ViT face & 57\% / 0.57 & 58\% / 0.70 \\
Lookup + VLM (60-epoch cap) & 61\% / 0.60 & 69\% / 0.75 \\
\best{Lookup + VLM (150-epoch cap)} & 62\% / 0.60 & 72\% / \best{0.77} \\
Lookup + VLM + uploaded ViT face & 58\% / 0.56 & 66\% / 0.72 \\
Lookup + VLM + our FER face (150-epoch cap) & 64\% / 0.69 & 69\% / 0.75 \\
Lookup + VLM + pose with speed (150-epoch cap) & 58\% / 0.59 & 59\% / 0.67 \\
\quad as above, VLM also fed to the head & 58\% / 0.59 & 60\% / 0.67 \\
\midrule
VLM only, no annotation (60-epoch cap) & 48\% / 0.53 & \best{76\%} / 0.69 \\
\bottomrule
\end{tabular}
\end{table}

\begin{table}[!htbp]
\centering
\caption{\textbf{Recurrent baselines (GRU, BiLSTM) vs ContextPose by input stream on the binary calls} (153 development clips, the same video-grouped 5-fold cross-validation, votes objective, 5-seed ensembles). Cells are balanced accuracy (\%) / AUC.
\emph{Architectures:} the recurrent baselines read the three windows in time order as one 96-step sequence (region one-hot, pose, encoded face times its quality mask, and the mask), project each step to width 64, and run a GRU (64) or a BiLSTM ($2 \times 32$); the final state is concatenated with the projected context vector and passed to an MLP head. ContextPose instead attends over each window with a context query and averages only the gated face steps.
\emph{Offsets:} rows with the annotation start from the event-annotation lookup; rows without it start from a constant offset. In the VLM rows the VLM embedding is the context vector and the annotation enters through the lookup.
\emph{Training:} recurrent runs up to 150 epochs (patience 10); ContextPose rows up to 60 epochs, except the three VLM rows (150). Annotation-only and annotation + VLM have no sequence for a recurrence to read (--). Valence: positive vs not (150 clips, 100 positive); arousal: high vs low (148 clips, 33 high). Bold: best per target.}
\label{tab:rnn}
\small
\begin{tabular}{lcccccc}
\toprule
 & \multicolumn{3}{c}{Valence} & \multicolumn{3}{c}{Arousal} \\
\cmidrule(lr){2-4}\cmidrule(lr){5-7}
Inputs & ContextPose & GRU & BiLSTM & ContextPose & GRU & BiLSTM \\
\midrule
Pose only & 46 / .52 & 45 / .51 & 52 / .45 & 53 / .57 & 59 / .43 & 57 / .40 \\
Annotation only & 55 / .58 & -- & -- & 68 / .61 & -- & -- \\
Face only & 60 / \best{.73} & 50 / .62 & 54 / .64 & 65 / .69 & 46 / .54 & 64 / .59 \\
\midrule
Pose + annotation & 54 / .60 & 57 / .59 & 57 / .57 & 62 / .68 & 62 / .64 & 60 / .62 \\
Pose + face & 50 / .58 & 62 / .65 & 57 / .59 & 55 / .61 & 56 / .50 & 57 / .55 \\
Annotation + face & \best{65} / .71 & 55 / .64 & 53 / .63 & 69 / .73 & 64 / .67 & 67 / .67 \\
\midrule
Pose + annotation + face & 58 / .63 & 60 / .61 & 60 / .63 & 60 / .69 & 67 / .64 & 56 / .68 \\
\midrule
Annotation + VLM & 62 / .60 & -- & -- & \best{72 / .77} & -- & -- \\
Pose + annotation + VLM & 58 / .59 & 62 / .60 & 54 / .58 & 59 / .67 & 66 / .70 & 63 / .68 \\
Annotation + VLM + face & 64 / .69 & 61 / .65 & 60 / .64 & 69 / .75 & 64 / .69 & 68 / .74 \\
\bottomrule
\end{tabular}
\end{table}

\end{document}
